\section{Formal Statements}
\label{sec:Formalization}

\begin{theorem}[Witness packaging]
Assume there exists a finitely generated group together with a finite generating alphabet such
that:
\begin{enumerate}[label=(\roman*)]
  \item the group has solvable word problem, and
  \item its geodesic growth is irrational.
\end{enumerate}
Then Br\"onnimann's Question~3 has a positive answer.
\end{theorem}

\begin{proof}[Formal content]
This is a direct unpacking of the existential definition of a positive answer.
\end{proof}

\begin{proposition}[Coordinate-model realization]
The virtually Engel presentation admits a canonical homomorphism into an explicit computable
semidirect product in which the generator $a$ is represented by a diagonal step in the nilpotent
factor and the generator $t$ by the non-trivial order-two symmetry.
\end{proposition}

\begin{proof}[Formal content]
The coordinate group law and the reflection automorphism are defined explicitly on integer
coordinates. One then checks directly that the four defining relators of Bodart's presentation are
satisfied in this target, so the universal property of presented groups yields the homomorphism.
\end{proof}

\begin{proposition}[Lattice constraints]
Inside the nilpotent coordinate factor, the subgroup generated by the diagonal steps is cut out by
two congruence conditions on the area and barycentric coordinates. These conditions are preserved
by multiplication, inversion, and the order-two symmetry.
\end{proposition}

\begin{proof}[Formal content]
The two residues are defined by explicit polynomial expressions in the coordinates. Their behavior
under multiplication, inversion, and reflection is computed directly, and divisibility by $2$ and
$6$ is propagated by elementary parity arguments.
\end{proof}

\begin{proposition}[Hall normal-form reconstruction]
Every point of the coordinate lattice is assigned an explicit Hall normal-form word whose image in
the coordinate model is exactly that lattice point.
\end{proposition}

\begin{proof}[Formal content]
The middle commutator in the normal form is taken in Hall's convention
$a^{-1}(a^t)^{-1}a a^t$, rather than the literal mathlib commutator $[a,a^t]$. With this choice,
the area and barycentric residues recover the exponents of the Hall commutators by exact division
by $2$ and $6$, and substituting these exponents back into the coordinate formula reproduces the
given lattice point.
\end{proof}

\subsection*{Word problem via injectivity of the coordinate map}

\begin{proposition}[Word-problem reduction]
Injectivity of the coordinate homomorphism suffices for solvable word problem: given
two finite words over the presentation alphabet, equality of the elements they represent in the
presented group is equivalent to equality of their images in the coordinate model, and the latter
is decidable because the coordinate model is a semidirect product of two types with decidable
equality.
\end{proposition}

\begin{proof}[Formal content]
The word-problem witness is stated as decidability of equality of pairs of finite words in the
presented group. Applying the coordinate homomorphism to both sides preserves equality; the
converse implication is exactly injectivity. On the coordinate side, equality is decided
componentwise on the underlying integer tuples and the two-element symmetry factor.
\end{proof}

\subsection*{Framework for the injectivity theorem}

The internal framework now includes a canonical set-theoretic retract $r\colon
\operatorname{CoordGroup} \to \operatorname{PresentedGroup} R$ whose value on a coordinate that
lies in the coordinate lattice is the equivalence class of its Hall normal-form word. The
following facts are proved internally:
\begin{itemize}
  \item every image of a word in the presented group lies in the coordinate lattice
    (\emph{lattice closure of the coordinate map});
  \item injectivity of the coordinate homomorphism is equivalent to $r$ being a left inverse of
    the coordinate homomorphism;
  \item on each of the presentation generators, $r$ composed with the coordinate homomorphism
    returns the same generator inside the presented group.
\end{itemize}
Together these reduce the remaining internal step to a single multiplication-compatibility
identity: for all pairs of elements of the presented group, the retract of the product of their
coordinate images equals the product of their retracts. The Hall commutation identities implied
by the four defining relators of the presentation are precisely what is needed to establish this
identity.

\subsection*{Remaining internal step}
The only remaining algebraic obligation is injectivity of the coordinate homomorphism, i.e.\ the
statement that two words representing distinct elements in the presented group are separated by
their integer coordinates. This is the Hall normal-form theorem inside the presented group: every
element is equal in the presented group to $\operatorname{mk}$ of the explicit Hall normal-form
word attached to its coordinates. The coordinate-model scaffolding, the lattice-congruence
subgroup, and the Hall normal-form reconstruction are all designed to feed this step. Once
injectivity is proved internally, the witness-packaging theorem can be instantiated with Bodart's
virtually Engel group and the remaining literature dependency reduces to Bodart's asymptotic
estimate.
